\section{Natural branch-local verification refuses blind repair (RQ4)}
\label{sec:rq4}

\begin{figure}[t]
\centering
\includegraphics[width=\linewidth]{fig4_branch_natural.pdf}
\caption{Direct branch-local audit of the FBCSP-LGG fusion backbone (Phase 4B refit, 21 leave-one-subject-out
folds on BNCI2014\_001/2015\_001). The spatial branch is simultaneously the most subject-leaky (L1 subject
decodability $0.92$/$0.82$), the most load-bearing (L4 ablation drop $+0.086$; highest fusion-gate weight), and
functionally coupled (L5 head-replay specificity). Yet \emph{removing} its subject subspace \emph{hurts} the
target (L6 task drop positive, $+0.050^{*}$ on 2015): the erasing-helps criterion fails. Verdict:
\textsc{no verified harmful branch shortcut}.}
\label{fig:branch}
\end{figure}

A multi-branch backbone lets us ask whether leakage \emph{means} something different in a load-bearing branch
than elsewhere --- and, decisively, whether the strongest natural leakage candidate is actually a harmful
shortcut. Unlike the frozen-artifact stage, we now measure the full branch-local ladder directly: a minimal
ERM re-fit of the same FBCSP-LGG recipe (no CMI, no architecture or hyperparameter search, no target-label fit)
reproduces the frozen deployment metrics (target bAcc $0.358$/$0.599$ vs.\ $0.349$/$0.608$, $|\Delta|{=}0.009$)
and, additionally, dumps per-branch latents and best-epoch checkpoints, so L1--L6 are computed per branch on real
EEG. The target-label firewall is clean on all $21$ folds (labels enter only final scoring; branch selection,
probes, and subspace fits are source-only).

\paragraph{The spatial branch is the strongest candidate on every low level, yet not harmful.} The spatial branch
is the most subject-decodable (L1 $0.92$/$0.82$), the most load-bearing (zeroing it drops task-bAcc by
$+0.086$ pooled; highest fusion-gate weight), and functionally coupled (L5: erasing its source-estimated subject
subspace changes the head's output specifically, beyond a random-subspace control). By L1--L5 alone it looks like
the canonical harmful shortcut. But the target consequence inverts the inference: removing the subject subspace
\emph{raises} the erased-model's task drop rather than lowering it (L6 $\text{task\_drop}{=}+0.015$ on 2a,
$+0.050^{*}$ on 2015), i.e.\ \emph{erasing the subject signal hurts the target}. The subject subspace is
task-entangled --- a task-useful reliance, not a harmful one.

\paragraph{Verification, not a leap.} FSR therefore \textbf{refuses} to call the strongest natural subject-leakage
branch a harmful shortcut: the verdict is \textsc{no verified harmful branch shortcut}
(\claim{C16}). This is the ladder working as designed --- a claim that stops at L1/L5 (``subject signal is
strongly decodable and functionally coupled, so remove it'') would be wrong here, because the L6 test shows the
removal is harmful. We can state that the spatial branch is load-bearing (\claim{C6}) and that its subject signal
is measurable and coupled, but we \emph{cannot} say ``spatial leakage is harmful'' or ``erase spatial subject
directions to improve DG.'' On natural EEG, the honest conclusion is a refusal.
